\documentclass[11pt]{article}
\usepackage[a4paper,margin=24mm]{geometry}
\usepackage{amsmath,amssymb}
\setlength{\parindent}{0pt}
\setlength{\parskip}{6pt}
\newcommand{\V}{V}
\newcommand{\E}{\mathrm e}
\begin{document}
\begin{center}
{\Large\bfseries Full budget permits zero oracle queries}\\[7pt]
{\small A counterexample to Conjecture 00000007863's linear lower bound}
\end{center}

\textbf{The source claim.} The bilingual source asserts that the expected
query complexity \(q(n,k)\) of randomized value-oracle maximization of a
normalized monotone submodular function is \(\Theta(n)\) for a
\((1-1/\E)\)-approximation, with exact lower-bound constant one and every
point queried at least once. It imposes no exclusion of \(k=n\), no
fixed-\(k\) asymptotic, and no restriction separating \(k\) from \(n\).
We disprove the necessary positive linear lower-bound conjunct on the
unbounded admissible family \(k=n\).

\textbf{Every objective is handled.} Let \(\V=\{0,\ldots,n-1\}\).
The allowed objective is an actual set function
\(f:2^{\V}\to\mathbb R\) satisfying
\[
 f(\varnothing)=0,\quad A\subseteq B\Longrightarrow f(A)\leq f(B),
 \quad f(A\cup B)+f(A\cap B)\leq f(A)+f(B).
\]
At budget \(k=n\), the ground set itself is feasible. Every feasible
set \(S\) is a subset of \(\V\), so monotonicity gives
\[
 f(S)\leq f(\V).
\]
Thus \(\V\) attains the actual maximum, for \emph{every} objective in
the source class. This does not select an easy instance such as the
zero function. The cardinality objective \(f(S)=|S|\), for example, is
normalized, monotone and modular, and has positive optimum \(n\) when
\(n>0\). It witnesses that the class and counterexample are nonvacuous.

\textbf{Approximation guarantee.} Normalization and monotonicity imply
\(0=f(\varnothing)\leq f(S)\) for every set. Put
\(\alpha=1-1/\exp(1)\), so \(0<\alpha\leq1\). Consequently
\[
 \alpha f(S)\leq f(S)\leq f(\V).
\]
Returning \(\V\) therefore satisfies the requested approximation,
indeed exact optimality, uniformly over the whole class. Submodularity
is retained in the hypotheses, even though monotonicity alone makes
this full-budget argument work.

\textbf{Actual zero-query protocol.} The algorithm is simply
\[
 \text{return }\V.
\]
It does not use an oracle answer or depend on which objective is
present. In the formal model, a protocol is built from return nodes
and query nodes; a query requests \(f(S)\), feeds that real answer to
its continuation, and increases the interpreter's query count by one.
The return node evaluates to output \(\V\) and count zero for every
oracle. Zero cost is proved from this actual interpreter, not supplied
as a premise or assigned to an unrelated complexity scalar.

\newpage
\begin{center}
{\large\bfseries Expected cost, asymptotics, and precise scope}
\end{center}

\textbf{A legitimate randomized special case.} A deterministic protocol
is allowed among randomized algorithms. Explicitly, take a singleton
random seed with probability mass one. The seed law is nonnegative
and normalized. Its expected query count is the finite weighted sum
of the actual interpreter counts. For the same return-ground-set
protocol this gives
\[
 \mathbb E[\text{queries}\mid f]=1\cdot0=0
 \qquad\text{for every allowed }f.
\]
Therefore even the worst case over all source objectives has zero
expected queries. Query counts are nonnegative, so this uniformly
valid algorithm already rules out any positive lower bound on the
least required worst-case expected count at full budget.

\textbf{Asymptotic contradiction.} Suppose a positive linear lower
bound held eventually along \((n,k)=(n,n)\). There would be \(c>0\)
and a threshold such that every valid randomized approximation
algorithm, for every sufficiently large \(n\), had expected cost at
least \(cn\) on some allowed objective. Apply that claim to the
return-ground-set protocol. Its expected cost is zero on every
objective, whereas \(cn>0\) once \(n\geq1\). This is a contradiction.
Thus no eventual positive \(\Omega(n)\) lower comparison is possible
on this unbounded family, and the stated unrestricted
\(\Theta(n)\) law, including its exact-constant assertion, fails.

\textbf{Scope of the conclusion.} We use the standard optimization
output: a feasible set with the promised objective guarantee. The
source does not request separately reporting its unknown objective
value. It also does not restrict the budget to \(k<n\) or fix \(k\)
as \(n\) grows. Those restrictions would describe a different claim.
Independent source-bound review must assess these conventions.

The source additionally asserts a query jump for improved
approximation and a uniquely determined information-theoretic
exponent. Those clauses are not individually formalized or settled
here. A false necessary linear lower-bound conjunct suffices to
refute the original conjunction on its stated domain. The simple
protocol class is used only to give an allowed counterexample;
it does not purport to rebuild the full randomized oracle model.

\textbf{Lean correspondence.} \texttt{SourceObjective} retains all
three source properties on actual \texttt{Finset (Fin n)} objects.
The complete source proves greatest-value optimality, a nontrivial
cardinality witness, and the exact approximation for every objective.
\texttt{ValueProgram} and \texttt{run} implement actual oracle queries
and their cost; the normalized seed expectation is proved zero.
The expected-query lower-bound theorem supplies the asymptotic
contradiction. The premise-free theorem
\texttt{full\_budget\_counterexample} collects the full certificate.

\textit{Provenance and verification.} This elementary full-budget route
was identified by the root coordinator and formalized in the same-run
scout before this author package. No first-discovery claim is made.
The packet is an author artifact; deterministic replay, independent
semantic review and acceptance belong to the manager.
\end{document}
